% Part:sets-functions-relations
% Chapter: size-of-sets
% Section: enumerations-alt

\documentclass[../../../include/open-logic-section]{subfiles}

\begin{document}

\olfileid{sfr}{siz}{enm-alt}

\olsection{তালিকায়ন ও \usetoken{S}{enumerable} সেট}

\begin{editorial}
  এই অংশে তালিকায়নকে $\Nat$-এর (প্রারম্ভিক খণ্ডগুলির) সঙ্গে একৈক সমাপতন হিসেবে সংজ্ঞায়িত করা হয়েছে, যেমনটি সেটতত্ত্বে করা হয়। তাই \olref[enm]{sec}-এর সংজ্ঞাগুলির সঙ্গে এর সামান্য অমিল আছে, এবং সেখানে থাকা সব উদাহরণ এখানে আবার দেওয়া হয়েছে। এই অংশটি ওই অংশের তুলনায় কিছুটা সংক্ষিপ্তও।
\end{editorial}

একটি সসীম সেটকে আমরা তার !!{element}s স্রেফ তালিকায় লিখেই নির্দিষ্ট করতে পারি। যেমন, কোনো সেটকে এভাবে সংজ্ঞায়িত করার সময় আমরা তা-ই করি:
\[
  A = \{a_1, a_2, \ldots, a_n\}.
\]
!!{element}s $a_1$, \dots, $a_n$ সকলেই পরস্পর ভিন্ন ধরে নিলে, এর ফলে $A$ এবং প্রথম $n$টি স্বাভাবিক সংখ্যা $0$, \dots, $n-1$-এর মধ্যে !!a{bijection} পাওয়া যায়। বিপরীতভাবে, প্রত্যেক সসীম সেটে কেবল সসীমসংখ্যক !!{element}s থাকে বলে, প্রত্যেক সসীম সেটকেই এমন একটি সমাপতনে আনা যায়। অন্যভাবে বললে, $A$ সসীম হলে $A$ এবং $\{0, \dots, n-1\}$-এর মধ্যে !!a{bijection} আছে, যেখানে $n$ হলো~$A$-এর !!{element}s-এর সংখ্যা।

নির্দিষ্ট ধরনের অসীম সেটও গ্রহণ করলে, কিছু অসীম সেটকেও তালিকায়িত হতে দেব। কোনো অসীম সেট এবং সমগ্র~$\Nat$-এর মধ্যে !!a{bijection} থাকলে সেটি তালিকায়িত হয়েছে—এ কথা বলে আমরা ধারণাটিকে নিখুঁত করতে পারি।

\begin{defn}[তালিকায়ন, সেটতাত্ত্বিকভাবে]
কোনো সেট $A$-এর একটি \emph{তালিকায়ন} হলো এমন !!a{bijection}, যার মানসমষ্টি $A$ এবং যার সংজ্ঞাক্ষেত্র হয় স্বাভাবিক সংখ্যার কোনো প্রারম্ভিক সেট $\{0, 1, \ldots, n\}$, {অথবা} স্বাভাবিক সংখ্যার সমগ্র সেট~$\Nat$।
\end{defn}

\begin{explain}
\emph{তালিকায়ন} শব্দটির এই ব্যবহারের পিছনে একটি স্বজ্ঞাত ধারণা আছে। কোনো সেট $A$-কে আমরা তালিকায়িত করেছি বলার অর্থ হলো, এমন !!a{bijection} $f$ আছে যার সাহায্যে সেট $A$-এর উপাদানগুলি একে একে গোনা যায়। $0$-তম উপাদান $f(0)$, প্রথমটি $f(1)$, \ldots, $n$-তমটি $f(n)$\ldots।\footnote{হ্যাঁ, আমরা $0$ থেকে গোনা শুরু করি। অবশ্যই~$1$ থেকেও শুরু করতে পারতাম। তাতে বড় কোনো পার্থক্য হতো না। শুধু~$\Nat$-এর বদলে~$\PosInt$ বসাতে হতো।} নিচের সংজ্ঞাটি যোগ করলে এর যুক্তি আরও স্পষ্ট হয়:
\end{explain}

\begin{defn}
  \ollabel{defn:enumerable}
  একটি সেট~$A$ !!{enumerable} হয় যদি এবং কেবল যদি হয় $A = \emptyset$, নয়তো~$A$-এর একটি তালিকায়ন থাকে। আমরা $A$-কে !!{nonenumerable} বলি যদি এবং কেবল যদি $A$ !!{enumerable} না হয়।
\end{defn}

\begin{explain}
অতএব, কোনো সেট !!{enumerable} হয় যদি এবং কেবল যদি সেটি শূন্য হয় অথবা কোনো তালিকায়নের সাহায্যে তার !!{element}s একে একে গোনা যায়।
\end{explain}

\begin{ex}
স্বাভাবিক সংখ্যাগুলিকে তালিকায়িত করে এমন অপেক্ষক হলো স্রেফ অভেদ অপেক্ষক $\Id{\Nat} \colon \Nat \to \Nat$, যেখানে $\Id{\Nat}(n) = n$। \emph{ধনাত্মক} স্বাভাবিক সংখ্যা $\Nat^+ = \Nat \setminus \{0\}$-কে তালিকায়িত করে $g(n) = n + 1$ অপেক্ষকটি, অর্থাৎ উত্তরসূরি অপেক্ষক।
\end{ex}

\begin{prob}
দেখাও যে কোনো সেট $A$ !!{enumerable} হয় যদি এবং কেবল যদি হয় $A = \emptyset$, নয়তো !!a{surjection} $f\colon \Nat \to A$ আছে। আরও দেখাও যে $A$ !!{enumerable} হয় যদি এবং কেবল যদি !!a{injection} $g\colon A \to \Nat$ আছে।
\end{prob}

\begin{ex}
নিচের সংজ্ঞা দ্বারা নির্ধারিত অপেক্ষক $f\colon \Nat \to \Nat$ ও $g \colon \Nat \to \Nat$
\begin{align*}
  f(n) & = 2n \text{ এবং}\\
  g(n) & = 2n+1
\end{align*}
যথাক্রমে জোড় স্বাভাবিক সংখ্যা এবং বিজোড় স্বাভাবিক সংখ্যাগুলিকে তালিকায়িত করে। কিন্তু কোনোটিই !!{surjective} নয়, তাই কোনোটিই $\Nat$-এর তালিকায়ন নয়।
\end{ex}

\begin{prob}
বর্গসংখ্যা $1$, $4$, $9$, $16$, \dots{}-এর একটি তালিকায়ন সংজ্ঞায়িত করো।
\end{prob}

\begin{ex}
$\lceil x \rceil$ দ্বারা \emph{ঊর্ধ্ব পূর্ণাংশ} অপেক্ষকটি বোঝানো হোক, যা $x$-কে উপরের দিকে নিকটতম পূর্ণসংখ্যায় নিয়ে যায়। তাহলে নিচের সংজ্ঞা দ্বারা নির্ধারিত অপেক্ষক $f \colon \Nat \to \Int$
\[
  f(n) = (-1)^{n} \left\lceil\tfrac{n}{2}\right\rceil
\]
পূর্ণসংখ্যার সেট~$\Int$-কে এভাবে তালিকায়িত করে:
\[
\begin{array}{c c c c c c c c}
f(0) & f(1) & f(2) & f(3) & f(4) & f(5) & f(6) & \dots \\ \\
\big\lceil \tfrac{0}{2} \big\rceil & -\big\lceil \tfrac{1}{2}\big\rceil &  \big\lceil \tfrac{2}{2} \big\rceil & -\big\lceil \tfrac{3}{2} \big\rceil & \big\lceil \tfrac{4}{2} \big\rceil  & -\big\lceil \tfrac{5}{2}\big\rceil & \big\lceil \tfrac{6}{2} \big\rceil & \dots \\ \\
0 & -1 & 1 & -2 & 2 & -3 & 3& \dots
\end{array}
\]
লক্ষ করো, ধনাত্মক ও ঋণাত্মক পূর্ণসংখ্যার মধ্যে এদিক-ওদিক ``লাফিয়ে'' $f$ কীভাবে $\Int$-এর মানগুলি তৈরি করে। নিচের মতো ক্ষেত্রভেদে সংজ্ঞায়িত অপেক্ষক হিসেবেও $f$-কে ভাবতে পারো:
\[
f(n) = \begin{cases}
  \frac{n}{2} & \text{যদি $n$ জোড় হয়}\\
  -\frac{n+1}{2} & \text{যদি $n$ বিজোড় হয়}
  \end{cases}
\]
\end{ex}

\begin{prob}
দেখাও যে $A$ ও $B$ !!{enumerable} হলে $A \cup B$-ও তাই।
\end{prob}

\begin{prob}
$n$-এর উপর গাণিতিক আরোহের সাহায্যে দেখাও যে $A_1$, $A_2$, \dots, $A_n$ সকলেই !!{enumerable} হলে $A_1 \cup \dots \cup A_n$-ও তাই।
\end{prob}

\end{document}
